% MathLab: acabamento editorial do livro (XeLaTeX). Incluído pelo Pandoc com -H.
% Fontes: TeX Gyre Pagella (texto), Fira Sans ou TeX Gyre Heros (títulos). Se faltarem, cai para Latin Modern.
\usepackage{xcolor}
\definecolor{ml-navy}{HTML}{14264F}
\definecolor{ml-ouro}{HTML}{F2B705}
\definecolor{ml-verde}{HTML}{22A559}
\definecolor{ml-laranja}{HTML}{F28A1F}
\definecolor{ml-roxo}{HTML}{7A4BD1}
\definecolor{ml-azul}{HTML}{1E6FD9}
\definecolor{ml-vermelho}{HTML}{D6453D}
\definecolor{ml-teal}{HTML}{0E8F8F}
\definecolor{ml-cinza}{HTML}{55607A}

\usepackage{etoolbox}
\providecommand{\mlsubtitulo}{}
% o Pandoc cola o subtítulo no título; aqui ele vai para a capa, em linha própria
\def\subtitle#1{\gdef\mlsubtitulo{#1}}
\IfFontExistsTF{TeX Gyre Pagella}{\setmainfont{TeX Gyre Pagella}}{}
\IfFontExistsTF{Fira Sans}{\setsansfont{Fira Sans}}{\IfFontExistsTF{TeX Gyre Heros}{\setsansfont{TeX Gyre Heros}}{}}
\IfFontExistsTF{TeX Gyre Pagella Math}{\setmathfont{TeX Gyre Pagella Math}}{}

\usepackage{enumitem}
\setlist{itemsep=0.15em,topsep=0.35em,parsep=0pt}
\setlist[itemize]{leftmargin=1.3em}
\setlist[enumerate]{leftmargin=1.7em}
\setlength{\emergencystretch}{3em}
\widowpenalty=10000 \clubpenalty=10000

% ---------- capítulos, seções, cabeçalho e rodapé ----------
\usepackage{titlesec}
\titleformat{\chapter}[display]
  {\sffamily\bfseries}
  {{\color{ml-ouro}\fontsize{58}{58}\selectfont\thechapter}\\[-0.2em]{\large\color{ml-cinza}\MakeUppercase{\chaptername}}}
  {0.2em}
  {\color{ml-navy}\huge\raggedright}
  [{\vspace{0.5em}\color{ml-ouro}\titlerule[2.2pt]}]
\titleformat{name=\chapter,numberless}[display]
  {\sffamily\bfseries}{}{0pt}
  {\color{ml-navy}\huge\raggedright}
  [{\vspace{0.5em}\color{ml-ouro}\titlerule[2.2pt]}]
\titlespacing*{\chapter}{0pt}{-1.5em}{2.2em}
\titleformat{\section}{\sffamily\bfseries\Large\color{ml-navy}}{\thesection}{0.6em}{}
\titleformat{\subsection}{\sffamily\bfseries\large\color{ml-navy}}{\thesubsection}{0.6em}{}

\usepackage{fancyhdr}
\renewcommand{\chaptermark}[1]{\markboth{#1}{}}
\pagestyle{fancy}
\fancyhf{}
\fancyhead[L]{\sffamily\small\color{ml-cinza}\nouppercase{\leftmark}}
\fancyfoot[C]{\sffamily\small\color{ml-navy}\thepage}
\renewcommand{\headrule}{\hbox to\headwidth{\color{ml-navy!35}\leaders\hrule height \headrulewidth\hfill}}
\renewcommand{\headrulewidth}{0.4pt}
\fancypagestyle{plain}{\fancyhf{}\fancyfoot[C]{\sffamily\small\color{ml-navy}\thepage}\renewcommand{\headrulewidth}{0pt}}

% ---------- capa ----------
\makeatletter
\renewcommand{\maketitle}{%
  \begin{titlepage}
  \pagecolor{ml-navy}\color{white}\thispagestyle{empty}
  \vspace*{0.14\textheight}
  {\noindent\color{ml-ouro}\rule{0.22\textwidth}{3pt}\par}
  \vspace{1.6em}
  \begin{flushleft}\setlength{\parindent}{0pt}
  {\sffamily\bfseries\fontsize{34}{40}\selectfont\@title\par}
  \ifdefempty{\mlsubtitulo}{}{\vspace{0.9em}{\sffamily\Large\color{ml-ouro}\mlsubtitulo\par}}
  \end{flushleft}
  \vfill
  \noindent{\sffamily\Large\bfseries\@author\par}
  \vspace{0.4em}
  \noindent{\sffamily\small\color{ml-ouro}\MakeUppercase{MathLab Aulas Visuais}\par}
  \vspace{2.2em}
  \end{titlepage}
  \clearpage\nopagecolor\color{black}%
}
\makeatother

% ---------- caixas didáticas ----------
\usepackage{tcolorbox}
\tcbuselibrary{skins,breakable}
\newcommand{\caixacab}[3]{\ifstrempty{#2}{}{{\sffamily\bfseries\footnotesize\color{#1}\MakeUppercase{#2}\par\nobreak}\vspace{1pt}}%
  \ifstrempty{#3}{}{{\sffamily\bfseries\color{black!85}#3\par\nobreak}\vspace{3pt}}}
\tcbset{caixabase/.style={enhanced,breakable,frame hidden,boxrule=0pt,arc=2pt,
  colback=#1!6,borderline west={3pt}{0pt}{#1},left=11pt,right=9pt,top=6pt,bottom=6pt,
  before skip=0.9em,after skip=0.9em,
  before upper={\setlength{\parindent}{0pt}\setlength{\parskip}{0.45em}}}}
\newtcolorbox{caixaobjetivos}[2]{caixabase=ml-navy,before upper={\caixacab{ml-navy}{#1}{#2}}}
\newtcolorbox{caixaprevios}[2]{caixabase=ml-cinza,before upper={\caixacab{ml-cinza}{#1}{#2}}}
\newtcolorbox{caixaproblematizacao}[2]{caixabase=ml-ouro,before upper={\caixacab{ml-ouro!70!black}{#1}{#2}}}
\newtcolorbox{caixacontexto}[2]{caixabase=ml-azul,before upper={\caixacab{ml-azul}{#1}{#2}}}
\newtcolorbox{caixaconceito}[2]{caixabase=ml-navy,before upper={\caixacab{ml-navy}{#1}{#2}}}
\newtcolorbox{caixaexemplo}[2]{caixabase=ml-verde,before upper={\caixacab{ml-verde!80!black}{#1}{#2}}}
\newtcolorbox{caixavisual}[2]{caixabase=ml-teal,before upper={\caixacab{ml-teal}{#1}{#2}}}
\newtcolorbox{caixaatividade}[2]{caixabase=ml-laranja,before upper={\caixacab{ml-laranja!85!black}{#1}{#2}}}
\newtcolorbox{caixaexercicios}[2]{caixabase=ml-roxo,before upper={\caixacab{ml-roxo}{#1}{#2}}}
\newtcolorbox{caixagabarito}[2]{caixabase=ml-azul,before upper={\caixacab{ml-azul}{#1}{#2}}}
\newtcolorbox{caixaavaliacao}[2]{caixabase=ml-roxo,before upper={\caixacab{ml-roxo}{#1}{#2}}}
\newtcolorbox{caixaacessibilidade}[2]{caixabase=ml-teal,before upper={\caixacab{ml-teal}{#1}{#2}}}
\newtcolorbox{caixanota}[2]{caixabase=ml-vermelho,before upper={\caixacab{ml-vermelho}{#1}{#2}}}
\newtcolorbox{caixadestaque}[2]{caixabase=ml-ouro,colback=ml-ouro!14,before upper={\caixacab{ml-ouro!60!black}{#1}{#2}}}

\usepackage{caption}
\captionsetup{font={small,sf},labelfont=bf,skip=4pt}
